%iffalse
\let\negmedspace\undefined
\let\negthickspace\undefined
\documentclass[journal,12pt,twocolumn]{IEEEtran}
\usepackage{cite}
\usepackage{amsmath,amssymb,amsfonts,amsthm}
\usepackage{algorithmic}
\usepackage{graphicx}
\usepackage{textcomp}
\usepackage{xcolor}
\usepackage{txfonts}
\usepackage{listings}
\usepackage{enumitem}
\usepackage{mathtools}
\usepackage{gensymb}
\usepackage{comment}
\usepackage[breaklinks=true]{hyperref}
\usepackage{tkz-euclide} 
\usepackage{listings}
\usepackage{gvv}                                        
%\def\inputGnumericTable{}                                 
\usepackage[latin1]{inputenc}                                
\usepackage{color}                                            
\usepackage{array}                                            
\usepackage{longtable}                                       
\usepackage{calc}                                             
\usepackage{multirow}                                         
\usepackage{hhline}                                           
\usepackage{ifthen}                                           
\usepackage{lscape}
\usepackage{tabularx}
\usepackage{array}
\usepackage{float}


\newtheorem{theorem}{Theorem}[section]
\newtheorem{problem}{Problem}
\newtheorem{proposition}{Proposition}[section]
\newtheorem{lemma}{Lemma}[section]
\newtheorem{corollary}[theorem]{Corollary}
\newtheorem{example}{Example}[section]
\newtheorem{definition}[problem]{Definition}
\newcommand{\BEQA}{\begin{eqnarray}}
\newcommand{\EEQA}{\end{eqnarray}}
\newcommand{\define}{\stackrel{\triangle}{=}}
\theoremstyle{remark}
\newtheorem{rem}{Remark}

% Marks the beginning of the document
\begin{document}
\bibliographystyle{IEEEtran}
\vspace{3cm}

\title{\textbf{Chapter-11 Limits, Continuity and Differentiability}}
\author{S. Sai Akshita EE24BTECH11054}
\maketitle
\newpage
\bigskip

\renewcommand{\thefigure}{\theenumi}
\renewcommand{\thetable}{\theenumi}
\begin{large}
\textbf{Section-B JEE Main/AIEEE}
\end{large}
\begin{enumerate}
    \item[(21)] Let $f$ be differentiable for all x.If $f$(1)=-2 and $f'(x)$ $\geq$ 2 for $x \in [1,6]$, then   \hfill[2005]
\begin{enumerate}
\item $f(6)$ \geq\)8 \item $f(6)<$8\item $f(6)<$5\item $f(6)$=5
\end{enumerate}
    \item[(22)] If $f$ is a real valued differentiable function satisfying $\mid f(x)-f(y)\mid \leq$(x-y)^2,$x,y\in R$ and $f(0)=0,$ then $f(1)$ equals \hfill[2005]
    \begin{enumerate}
    \item -1\item 0\item 2\item 1
    \end{enumerate}
    \item[(23)] Let $f$:R\rightarrow\) R be a function defined by\\$f(x)$=min\{x+1,\mid\)x\mid\)+1\}, Then which of the following is true? \hfill[2007]
    \begin{enumerate}
    \item $f(x)$ is differentiable everywhere\item $f(x)$ is not differentiable at x=0\item $f(x)$ \geq\)1 for all x\in\) R\item$f(x)$ is not differentiable at x=0
    \end{enumerate}
    \item[(24)] The function $f:R-\{0\}\rightarrow R$ given by\\$f(x)=\frac{1}{x}-\frac{2}{e^{2x}-1}$ can be made continuous at x=0 b defining $f(0)$ as \hfill[2007]
    \begin{enumerate}
        \item 0\item1\item 2\item -1
    \end{enumerate}
    \item[(25)] Let $f(x)= \left\{ \begin{array}{rcl}
(x-1)sin\frac{1}{x-1} & \mbox{if}
& x\neq 1 \\ 0 & \mbox{if} & x=1
\end{array}\right. $ \\Then which of the following is true? \hfill[2008]
    \begin{enumerate}
        \item f is neither differentiable at x=0 nor at x=1\item f is differentiable at x=0 and x=1\item f is differentiable at x=0 but not at x=1\item f is differentiable at x=1 but not at x=0
    \end{enumerate}
    \item[(26)] Let $f:R\rightarrow R$ be a positive increasing function with
    $\displaystyle \lim_{x\to\infty} \frac{f(3x)}{f(x)}=1 $.Then $\displaystyle \lim_{x\to\infty} \frac{f(2x)}{f(x)}$=\hfill[2010]
    \begin{enumerate}
     \item $\frac{2}{3}$ \item $\frac{3}{2}$ \item3 \item1
    \end{enumerate}
     \item[(27)]  $\displaystyle \lim_{x\to 2} \frac{\sqrt{1-cos{2(x-2)}}}{x-2} $ \hfill[2011]
    \begin{enumerate}
        \item equals $\sqrt{2}$ \item equals $-\sqrt{2}$ \item equals $1/\sqrt{2}$ \item does not exist
    \end{enumerate}
    \item[(28)] The values of p and q for which the function $f(x)= \left\{ \begin{array}{rcl}
\frac{sin(p+1)x+sinx}{x} & \mbox{,}
& x<0 \\ q & \mbox{,} & x=0 \\
 \frac{\sqrt{x+x^2}-\sqrt{x}}{x^{3/2}}& \mbox{,} & x>0
\end{array}\right. $ \\is  continuous for all x in R,are \hfill[2011]
    \begin{enumerate}
        \item p=5/2,q=1/2 \item p=-3/2,q=1/2 \item p=1/2,q=3/2 \item p=1/2,q=-3/2
    \end{enumerate}
\item[(29)] Let $f:R \rightarrow [0,\infty)$ be such that $\displaystyle \lim_{x\to 5} f(x)$ exists and  $\displaystyle \lim_{x\to 5}$ $\frac{(f(x))^2-9}{\sqrt{\mid x-5\mid}}$.Then $\displaystyle \lim_{x\to 5} f(x)$ equals:
\begin{enumerate}
    \item 0\item1\item2\item3
\end{enumerate}
\item[(30)] If $f:R\rightarrow\ R$ is a function defined by $f(x)$=[x]cos ($\frac{2x-1}{2}$) $\pi$, where [x] denotes greatest integer function,then $f$ is \hfill[2012]
\begin{enumerate}
    \item continuous for every real x \item discontinuous only at x=0 \item discontinuous only at non-zero integral values of x \item continuous only at x=0
    \end{enumerate}
    \item{(31)} Consider the function $f(x)=\mid x-2\mid +\mid x-5\mid$, x$\in$R.\\ \textbf{Statement-1:} $f'(4)=0$\\ \textbf{Statement-2:} $f$ is continuous in [2,5],differentiable in (2,5) and $f(2)=f(5)$ \hfill[2012]
    \begin{enumerate}
        \item Statement-1 is false, Statement-2 is true \item Statement-1 is true, Statement-2 is true;Statement-2 is a correct explanation of Statement-1 \item Statement-1 is true, Statement-2 is true; Statement-2 is \textbf{not} the correct explanation of Statement-1 \item Statement-1 is true, Statement-2 is false.
    \end{enumerate}
    \item[(32)] $\displaystyle \lim_{x \to 0}$  $\frac{(1-cos2x)(3+cosx)}{xtan4x}$ is equal to \hfill[\textbf{JEEM 2013}]
    \begin{enumerate}
        \item -$\frac{1}{4}$ \item $\frac{1}{2}$  \item1 \item2 \end{enumerate}
    \item[(33)] $\displaystyle \lim_{x \to 0}$ $\frac{sin(\pi\cos^{2}x)}{x^2}$ is equal to: \hfill[\textbf{JEEM 2014}]
        \begin{enumerate}
        \item $-\pi$ \item $\pi$ \item $\frac{\pi}{2}$ \item 1 
    \end{enumerate}
    \item[(34)] $\displaystyle \lim_{x \to 0}$$ \frac{(1-cos2x)(3+cosx)}{xtan4x}$is equal to \hfill[\textbf{JEEM 2015}]
        \begin{enumerate}
            \item 2 \item 1/2 \item4 \item3
        \end{enumerate}
        \item[(35)] If the function\\g(x)= \left\{ \begin{array}{rcl}
k\sqrt{x+1} & \mbox{,}
& 0\leq x\leq3 \\ mx+2 & \mbox{,} & 3<x\leq5
\end{array}\right. \\$ is differentiable, then the value of \\k+m is : \hfill[\textbf{JEEM 2015}]
        \begin{enumerate}  
            \item 10/3 \item 4 \item 2 \item 16/5
        \end{enumerate}


\end{document}
